
\medskip
\noindent\textbf{1.3.}\footnote{The editors expanded the original in this
paragraph.}
\phantomsection
\label{sga3:VIIA:1.3}
Let \(u:Y\to X\) be a morphism of \(S\)-schemes. Recall that the
adjunction isomorphism
\[
 \operatorname{Hom}_{p_X^{-1}(\mathcal O_S)}
   \bigl(\mathcal O_X,u_*(\mathcal O_Y)\bigr)
 \xrightarrow{\sim}
 \operatorname{Hom}_{p_Y^{-1}(\mathcal O_S)}
   \bigl(u^{-1}\mathcal O_X,\mathcal O_Y\bigr)
\]
sends a morphism of \(p^{-1}(\mathcal O_S)\)-modules
\[
  d:\mathcal O_X\longrightarrow u_*(\mathcal O_Y)
\]
to \(d'=\varepsilon\circ u^{-1}(d)\), where
\(\varepsilon:u^{-1}u_*(\mathcal O_Y)\to\mathcal O_Y\) is the
canonical morphism.

Write \(\mathcal J_u\) (respectively, \(\mathcal I_u\)) for the kernel
of the homomorphism of algebras
\[
 u^\natural:\mathcal O_X\longrightarrow u_*(\mathcal O_Y)
 \quad
 \left(\text{respectively, }
 u^{\natural\prime}:u^{-1}\mathcal O_X\longrightarrow\mathcal O_Y
 \right),
\]
and let \(d:\mathcal O_X\to u_*(\mathcal O_Y)\) be a morphism of
\(p^{-1}(\mathcal O_S)\)-modules. If \(U\subset X\) is open and
\(f_0,\ldots,f_n,g\in\mathcal O_X(U)\), induction on \(n\) shows that
condition \((\ast_n)\) is equivalent to
\begin{equation}
  0=
  \sum_{I\subset\{0,\ldots,n\}}
  (-1)^{|I|}
  u^\natural\bigl(f_{\{0,\ldots,n\}\setminus I}\bigr)d(f_Ig),
  \tag{\(\ast\ast_n\)}
  \label{sga3:VIIA:eq:starstar-n}
\end{equation}
where \(f_I\) denotes the product of the \(f_i\) for \(i\in I\).
Consequently, if \(d\) satisfies \((\ast_n)\), then \(d\) vanishes on
the ideal \(\mathcal J_u^{\,n+1}\).

Now suppose that \(Y=S\). Then \(u:S\to X\) is a section of
\(p:X\to S\), hence an immersion. On the one hand,
\[
 \varepsilon:u^{-1}u_*\mathcal O_S\longrightarrow\mathcal O_S
\]
is an isomorphism, and therefore \(u^{-1}\mathcal J_u=\mathcal I_u\).
On the other hand, there is an isomorphism
\begin{equation}
  u^{-1}\mathcal O_X\simeq\mathcal O_S\oplus\mathcal I_u.
  \tag{\(\star\)}
  \label{sga3:VIIA:eq:section-splitting}
\end{equation}
Suppose that \(d\) vanishes on \(\mathcal J_u^{\,n+1}\). Then
\(d'=\varepsilon\circ u^{-1}(d)\) vanishes on
\(\mathcal I_u^{\,n+1}\), and therefore satisfies the analogues
\((\ast\ast'_n)\) and \((\ast'_n)\) of
\((\ast\ast_n)\) and \((\ast_n)\) when
\[
 f_0,\ldots,f_n\in\mathcal I_u\bigl(u^{-1}(U)\bigr).
\]
Moreover,
\[
  (\operatorname{ad}a)(\varphi)=0
\]
for every \(a\in\mathcal O_S(u^{-1}U)\) and every morphism of
\(\mathcal O_{u^{-1}U}\)-modules
\(\varphi:u^{-1}\mathcal O_U\to\mathcal O_{u^{-1}U}\).
The splitting \eqref{sga3:VIIA:eq:section-splitting} therefore implies
that \(d'\) satisfies \((\ast'_n)\), and hence that \(d\) satisfies
\((\ast_n)\).

\medskip
\noindent\textbf{Lemma.}
\phantomsection
\label{sga3:VIIA:1.3-section-lemma}
\textit{If \(u:S\to X\) is a section of \(p:X\to S\), then \(d\) is an
\(S\)-deviation of \(u\) of order \(\leq n\) if and only if \(d'\)
vanishes on \(\mathcal I_u^{\,n+1}\).}

This interpretation generalizes as follows. Let \(u:Y\to X\) be any
\(S\)-morphism, and let \(\Gamma_u:Y\to Y\times X\) be its graph, with
components \(\operatorname{id}_Y\) and \(u\). Every \(S\)-deviation
\(d\) of \(u\) of order \(\leq n\) gives, by composition,
\[
  Y\xrightarrow{\mathrm{diag}}Y\times Y
   \mathop{\longrightarrow}^{d_Y}_{u_Y}Y\times X,
\]
a \(Y\)-deviation of \(\Gamma_u\) of order \(\leq n\), denoted
\(\Gamma_d\), the graph of \(d\).

Conversely, a \(Y\)-deviation \(e\) of \(\Gamma_u\) gives the composite
\(S\)-deviation
\[
 e_X=\operatorname{pr}_2\circ e:
 \qquad
 Y\mathop{\longrightarrow}^{e}_{\Gamma_u}Y\times X
  \xrightarrow{\operatorname{pr}_2}X.
\]
One has immediately \((\Gamma_d)_X=d\). The identity
\(\Gamma_{e_X}=e\) follows from the \(\mathcal O_Y\)-linearity of
\(e\).
Thus there is an isomorphism of \(\mathcal O_Y(Y)\)-modules
\[
\left\{
\begin{array}{c}
S\text{-deviations of }u\\[-2pt]
\text{of order }\leq n
\end{array}
\right\}
\xrightarrow{\sim}
\left\{
\begin{array}{c}
Y\text{-deviations of }\Gamma_u\\[-2pt]
\text{of order }\leq n
\end{array}
\right\},
\qquad d\longmapsto\Gamma_d.
\]
Moreover, \(d\) is an \(S\)-derivation of \(u\) if and only if
\(\Gamma_d\) is a \(Y\)-derivation of \(\Gamma_u\).

Let \(\mathcal I_{\Gamma_u}\) denote the kernel of the homomorphism of
algebras
\[
  \Gamma_u^{-1}\mathcal O_{Y\times X}\longrightarrow\mathcal O_Y
\]
corresponding to \(\Gamma_u\). The preceding lemma yields:

\medskip
\noindent\textbf{Proposition.}
\phantomsection
\label{sga3:VIIA:1.3-graph-proposition}
\textit{Let \(u:Y\to X\) be an \(S\)-morphism and
\(\Gamma_u:Y\to Y\times X\) its graph. The \(S\)-deviations of \(u\)
of order \(\leq n\) identify with the \(Y\)-deviations of
\(\Gamma_u\) of order \(\leq n\), and these are in bijection with}
\[
 \operatorname{Hom}_{\mathcal O_Y}
 \left(
  \frac{\Gamma_u^{-1}\mathcal O_{Y\times X}}
       {\mathcal I_{\Gamma_u}^{\,n+1}},
  \mathcal O_Y
 \right).
\]

\medskip
\noindent\textbf{1.3.1.}\footnote{This paragraph was added by the
editors.}
\phantomsection
\label{sga3:VIIA:1.3.1}
Return to the case in which \(u:S\to X\) is a section of \(p:X\to S\).
The homomorphism
\[
 \varphi:u^{-1}\mathcal O_X\longrightarrow\mathcal O_S
\]
has a section, denoted simply by \(g\mapsto g\cdot1\). Thus, with the
notation of \SGARefLink{sga3:VIIA:1.3}{1.3},
\[
  u^{-1}\mathcal O_X\simeq\mathcal O_S\oplus\mathcal I_u,
\]
and, for every section \(f\) of \(u^{-1}\mathcal O_X\),
\[
  f-\varphi(f)\cdot1\in\mathcal I_u.
\]

Let \(d\) be an \(S\)-deviation of \(u\) of order \(\leq1\), and let
\(d':u^{-1}\mathcal O_X\to\mathcal O_S\) be the
\(\mathcal O_S\)-morphism corresponding to \(d\). If \(a,b\) are
sections of \(u^{-1}\mathcal O_X\), then
\begin{align*}
0
 &=d'\bigl((a-\varphi(a)\cdot1)(b-\varphi(b)\cdot1)\bigr)\\
 &=d'(ab)-\varphi(a)d'(b)-\varphi(b)d'(a)+\varphi(ab)d'(1).
\end{align*}
Consequently, \(d\) is an \(S\)-derivation of \(u\) if and only if
\(d'(1)=0\).

\medskip
\noindent\textbf{Lemma.}
\phantomsection
\label{sga3:VIIA:1.3.1-derivation-lemma}
\textit{The \(S\)-derivations of \(u\) are precisely the
\(S\)-deviations of \(u\) of order \(1\) which vanish on the unit
section of \(\mathcal O_X\). They correspond to the
\(\mathcal O_S(S)\)-module}
\[
  \operatorname{Hom}_{\mathcal O_S}
  \bigl(\mathcal I_u/\mathcal I_u^2,\mathcal O_S\bigr),
\]
\textit{and there is an isomorphism of \(\mathcal O_S(S)\)-modules}
\[
 \operatorname{Dev}_{\leq1}(u)
 \simeq
 \mathcal O_S(S)\oplus\operatorname{Der}_S(u).
\]

Returning to the general case gives:

\medskip
\noindent\textbf{Corollary.}
\phantomsection
\label{sga3:VIIA:1.3.1-graph-corollary}
\textit{Let \(u:Y\to X\) be an \(S\)-morphism and
\(\Gamma_u:Y\to Y\times X\) its graph. There is a canonical
isomorphism of \(\mathcal O_Y(Y)\)-modules}
\[
 \operatorname{Der}_S(u)
 \simeq
 \operatorname{Der}_Y(\Gamma_u)
 \simeq
 \operatorname{Hom}_{\mathcal O_Y}
 \bigl(\mathcal I_{\Gamma_u}/\mathcal I_{\Gamma_u}^2,\mathcal O_Y\bigr).
\]

\medskip
\noindent\textbf{Definition 1.4.}
\phantomsection
\label{sga3:VIIA:1.4}
\textit{Let \(X\) be an \(S\)-scheme. An \(S\)-differential operator
(respectively, an \(S\)-differential operator of order \(\leq n\)) on
\(X\) is an \(S\)-deviation (respectively, an \(S\)-deviation of order
\(\leq n\)) of the identity morphism of \(X\).}

By \SGARefLink{sga3:VIIA:1.1}{1.1}, an \(S\)-differential operator of order \(\leq n\) is therefore
an endomorphism of the \(p^{-1}(\mathcal O_S)\)-module
\(\mathcal O_X\) satisfying the identities \((\ast_n)\) of \SGARefLink{sga3:VIIA:1.1}{1.1}. We
write \(\operatorname{Dif}^n_{X/S}\) for the
\(\Gamma(\mathcal O_S)\)-module
of \(S\)-differential operators of order \(\leq n\), and
\(\operatorname{Dif}_{X/S}\) for the module of all
\(S\)-differential operators.

As observed in \SGARefLink{sga3:VIIA:1.2}{1.2}, \(S\)-deviations of \(\operatorname{id}_X\) can be
composed. This makes \(\operatorname{Dif}_{X/S}\) a
\(\Gamma(\mathcal O_S)\)-algebra, called the \emph{algebra of
differential operators of \(X/S\)}.

Similarly, for every open subset \(V\subset X\), set
\[
 \mathcal{Dif}_{X/S}(V)
 =\operatorname{Dif}_{V/S}
 =\operatorname{Dev}(\operatorname{id}_V).
\]
By \SGARefLink{sga3:VIIA:1.1.3}{1.1.3} this defines a sheaf of \(\mathcal O_X\)-modules, called the
\emph{sheaf of \(S\)-differential operators on \(X\)}.

\medskip
\noindent\textbf{1.4.1.}
\phantomsection
\label{sga3:VIIA:1.4.1}
As in \SGARefLink{sga3:VIIB:1.3.prop}{1.3}, differential operators on \(X/S\) may be interpreted using
the graph of the identity morphism of \(X\), namely the diagonal
\[
  \Delta=\Delta_{X/S}:X\longrightarrow X\times X.
\]
Endow \(\mathcal O_{X\times X}\) with the
\(\operatorname{pr}_1^{-1}(\mathcal O_X)\)-algebra structure defined
by \(\operatorname{pr}_1\). Then
\(\Delta^{-1}\mathcal O_{X\times X}\) is an algebra over
\[
  \mathcal O_X=\Delta^{-1}\operatorname{pr}_1^{-1}\mathcal O_X.
\]
Let \(\mathcal I_{X/S}\) be the kernel of the homomorphism
\[
 \Delta^{-1}\mathcal O_{X\times X}
 \xrightarrow{m}
 \mathcal O_X
\]
adjoint to
\(\mathcal O_{X\times X}\to\Delta_*(\mathcal O_X)\), and set
\[
  \mathcal P^m_{X/S}
  =
  \Delta^{-1}\mathcal O_{X\times X}/\mathcal I_{X/S}^{\,m+1}.
\]
If \(V\subset S\) and \(U\subset X\) are affine open subsets with
\(U\) above \(V\), and if
\[
  k=\Gamma(V,\mathcal O_S),
  \qquad
  A=\Gamma(U,\mathcal O_X),
\]
then
\[
  \Gamma(U,\mathcal P^m_{X/S})
  =(A\otimes_kA)/I^{m+1},
\]
where \(I\) is generated by the elements \(a\otimes1-1\otimes a\),
for \(a\in A\).

By \SGARefLink{sga3:VIIB:1.3.prop}{1.3} there is an isomorphism of \(\mathcal O_X(X)\)-modules
\[
 j_X:\operatorname{Dif}^m_{X/S}
 \xrightarrow{\sim}
 \operatorname{Hom}_{\mathcal O_X}
   (\mathcal P^m_{X/S},\mathcal O_X).
\]
It is given explicitly as follows: if
\(d\in\operatorname{Dif}^m_{X/S}\) and a section \(c\) of
\(\mathcal P^m_{X/S}\) over \(U\) has the form
\[
  c=a\otimes b+I^{m+1},
\]
then \(j_X(d)(c)=a\,d(b)\).

\medskip
\noindent\textbf{1.4.2.}
\phantomsection
\label{sga3:VIIA:1.4.2}
Let \(d\) be a differential operator and \(u\) a section of \(X\) over
\(S\). The \emph{value of \(d\) at \(u\)} is the composite
\(S\)-deviation
\[
  S\xrightarrow{u}X
  \mathop{\longrightarrow}^{d}_{\operatorname{id}_X}X.
\]
By \SGARefLink{sga3:VIIB:1.3.prop}{1.3} and \SGARefLink{sga3:VIIA:1.4.1}{1.4.1}, if \(d\) has order \(\leq m\), then \(du\)
(respectively, \(d\)) is canonically associated with a morphism of
\(\mathcal O_S\)-modules
\[
 d':u^{-1}\mathcal O_X/\mathcal I_u^{\,m+1}\longrightarrow\mathcal O_S
\]
(respectively, a morphism of \(\mathcal O_X\)-modules
\(d'':\mathcal P^m_{X/S}\to\mathcal O_X\)).

The morphism \(d'\) may be recovered from \(d''\) as follows. The square
\begin{center}
  \phantomsection
  \label{sga3:VIIA:1.4.2-value-cartesian-square}
  \begin{tikzcd}[column sep=huge,row sep=large]
X\simeq S\times X
  \arrow[r,"u\times X"]
  \arrow[d,"p"']
&
X\times X
  \arrow[d,"\operatorname{pr}_1"]
\\
S
  \arrow[r,"u"]
&
X
\end{tikzcd}

\end{center}
is cartesian. Hence \(X\) identifies with
\(S\times_X(X\times X)\), \(u\) with \(S\times_X\Delta\), and
\(u^*(\mathcal P^m_{X/S})\) with
\(u^{-1}\mathcal O_X/\mathcal I_u^{\,m+1}\). Under these
identifications, \(u^*(d'')\) is the morphism
\[
  u^{-1}\mathcal O_X/\mathcal I_u^{\,m+1}\longrightarrow\mathcal O_S
\]
which is precisely \(d'\).

\medskip
\noindent\textbf{1.5.}
\phantomsection
\label{sga3:VIIA:1.5}
Set, as usual,
\[
  I_S=\operatorname{Spec}\mathcal O_S[T]/(T^2).
\]
Let \(\tau:S\to I_S\) be the zero section, and let \(\sigma\) be the
canonical deviation of \(\tau\) defined in \SGARefLink{sga3:VIIA:1.2.0}{1.2.0}: it is the
\(\mathcal O_S\)-module homomorphism which vanishes on the unit section
of \(\mathcal O_S[T]/(T^2)\) and sends the class \(t\) of \(T\) modulo
\(T^2\) to the unit section of \(\mathcal O_S\).

Let \(X\) be an \(S\)-scheme. Every \(I_S\)-automorphism \(u\) of
\(I_S\times X\) inducing the identity on \(X\) determines, by
composition, a differential operator \(D_u\) on \(X\):
\begin{center}
  \phantomsection
  \label{sga3:VIIA:1.5-infinitesimal-automorphism-sequence}
  \begin{tikzcd}[column sep=large]
X\simeq S\times X
  \arrow[r,"\sigma\times X"]
&
I_S\times X
  \arrow[r,"u"]
&
I_S\times X
  \arrow[r,"\operatorname{pr}_2"]
&
X\,.
\end{tikzcd}

\end{center}
By \SGARefLink{sga3:ii:subsection:3-14}{Expos\'e II, 3.14}, the map \(u\mapsto D_u\) is an isomorphism from
the \(\Gamma(\mathcal O_S)\)-Lie algebra
\[
  \operatorname{Lie}(\operatorname{Aut}X)
  :=\operatorname{Lie}(\operatorname{Aut}X)(S)
\]
onto the \(\Gamma(\mathcal O_S)\)-Lie algebra of
\(p^{-1}(\mathcal O_S)\)-derivations of \(\mathcal O_X\). Its inverse
sends a derivation \(D\) to the automorphism of \(I_S\times X\)
corresponding to
\[
  a+bt\longmapsto a+(Da+b)t
\]
on \(\mathcal O_X[T]/(T^2)\).

